% === Begin Label Data ===
% ID: 810
% 难度星级: 
% 标签: 
% 备注: 
% 组卷引用次数: 0
% === End  Label Data ===

\begin{problem}{2010}{G}{江西卷}{11}{概率}
一位国王的铸币大臣在每箱 $100$ 枚的硬币中各掺入了一枚劣币，国王怀疑大臣作弊，他用两种方法来检测．方法一：在 $10$ 箱中各任意抽查一枚；方法二：在 $5$ 箱中各任意抽查两枚．国王用方法一、二能发现至少一枚劣币的概率分别记为 $p_1$ 和 $p_2$，则 (\hspace{1cm})
\begin{choices}
\choice{{$p_1 = p_2$}}
\choice{{$p_1 < p_2$}}
\choice{{$p_1 > p_2$}}
\choice{{以上三种情况都有可能}}
\end{choices}
\end{problem}

\begin{answer}
B
\end{answer}

\begin{solutions}
在每箱中抽查一枚，找到劣币的概率是 $p = 0.01$．而在每箱中抽查两枚，找到劣币的概率是 $\dfrac{\mathrm{C}_{99}^{1}}{\mathrm{C}_{100}^{2}} = 0.02 = 2p$．  
所以，方法一能找到至少一枚劣币的概率是 $p_1 = 1 - (1 - p)^{10}$．方法二能找到至少一枚劣币的概率是 $p_2 = 1 - (1 - 2p)^5$．因为  
$$
p_1 = 1 - (1 - p)^{10} = 1 - (1 - 2p + p^2)^5 < 1 - (1 - 2p)^5 = p_2,
$$  
所以 $p_1 < p_2$．
\end{solutions}